\documentclass[11pt]{article}

\usepackage[margin=1in]{geometry}
\usepackage[T1]{fontenc}
\usepackage[utf8]{inputenc}
\usepackage{lmodern}
\usepackage{microtype}
\usepackage{amsmath,amssymb,amsthm,mathtools}
\usepackage{booktabs}
\usepackage{enumitem}
\usepackage{xcolor}
\usepackage{hyperref}
\usepackage[nameinlink,noabbrev]{cleveref}

\newcommand{\Prob}{\mathbb P}
\newcommand{\E}{\mathbb E}
\newcommand{\dd}{\mathrm d}
\newcommand{\e}{\mathrm e}
\newcommand{\R}{\mathsf R}
\newcommand{\Xrel}{X_{\mathrm{rel}}}
\newcommand{\Yrel}{Y_{\mathrm{rel}}}

\theoremstyle{plain}
\newtheorem{proposition}{Proposition}
\newtheorem{corollary}{Corollary}

\theoremstyle{definition}
\newtheorem{definition}{Definition}

\theoremstyle{remark}
\newtheorem{remark}{Remark}

\hypersetup{
  colorlinks=true,
  linkcolor=blue!55!black,
  citecolor=red!55!black,
  urlcolor=blue!60!black,
  pdftitle={Release counts from the killed generating function},
  pdfauthor={}
}

\title{Ultimate release counts from the killed generating function}
\author{}
\date{}

\begin{document}
\maketitle
\vspace{-1.5em}

\begin{abstract}
\noindent
For the two-type birth--death--conversion process with type-specific absorbing catastrophe, let
$(\Xrel,\Yrel)$ be the intracellular counts released at catastrophe, with the convention
$(\Xrel,\Yrel)=(0,0)$ when catastrophe never occurs.
We record their joint law, joint probability generating function, and marginals in terms of the
killed (defective) generating functions $\mathcal F_i$ and the catastrophe-flux generators
$\mathcal J_i$ from the catastrophe-aware PGF appendix.
The construction is exact for any non-negative type-specific rates; in the generic regime
$\delta_1,\delta_2>0$ the atom at the origin is ordinary extinction before absorption.
\end{abstract}

\section{Setup and definition}

We work throughout with the continuous-time two-type birth--death--conversion process of the main
paper: type~$1$ replicates at rate $\lambda_1$, dies at rate $\mu_1$, and converts one-way into
type~$2$ at rate $\nu$; type~$2$ replicates at rate $\lambda_2$ and dies at rate $\mu_2$.
In state $(x,y)$ a single global catastrophe occurs at rate $\delta_1x+\delta_2y$ and sends the
process to the absorbing cemetery $(\R,\R)$. Write $\tau_c$ for the catastrophe time
(with the usual convention $\tau_c=\infty$ on paths that never absorb).

Let $i\in\{1,2\}$ label the initial type, and write $\Prob_i$, $\E_i$ for the law and expectation
started from one individual of that type (so $i=1$ means $(X_0,Y_0)=(1,0)$ and $i=2$ means
$(0,1)$).

\begin{definition}[Ultimate release counts]
\label{def:release}
The ultimate type-$1$ and type-$2$ release counts are the random variables
\begin{equation}
\label{eq:release-def}
(\Xrel,\Yrel)
\;=\;
\begin{cases}
\bigl(X_{\tau_c-},\,Y_{\tau_c-}\bigr)
& \text{on }\{\tau_c<\infty\},\\[0.35em]
(0,0)
& \text{on }\{\tau_c=\infty\}.
\end{cases}
\end{equation}
Thus release records the left-limit population at the absorbing event when that event occurs, and
is set to zero when no release event ever takes place---in particular when the intracellular
population dies out by ordinary deaths before catastrophe.
\end{definition}

\begin{remark}[Generic regime versus boundaries]
\label{rem:generic}
In the generic regime $\delta_1>0$, $\delta_2>0$, the only way to have $\tau_c=\infty$ is ordinary
extinction to $(0,0)$ before the catastrophe clock fires. Consequently the atom of
$(\Xrel,\Yrel)$ at the origin is exactly the forever-avoidance probability, and coincides with
clearance without release. On certain catastrophe-rate boundaries (for instance $\delta_1=0$),
paths with $\tau_c=\infty$ need not be extinct; \cref{eq:release-def} still assigns $(0,0)$ on
those paths, now meaning ``no release event'' rather than ``empty burst''.
All formulae below remain valid on those boundaries; only the biological reading of the origin atom
changes.
\end{remark}

\section{Killed generating function and catastrophe flux}

Recall the killed (defective) generating functions
\begin{subequations}
\label{eq:killed-pgf}
\begin{align}
\mathcal F_1(x,y,t)
&=\E_{(1,0)}
\!\left[x^{X_t}y^{Y_t}\boldsymbol 1_{\{\tau_c>t\}}\right],
\\
\mathcal F_2(x,y,t)
&=\E_{(0,1)}
\!\left[x^{X_t}y^{Y_t}\boldsymbol 1_{\{\tau_c>t\}}\right],
\end{align}
\end{subequations}
equivalently the occupation-time transforms
\begin{equation*}
\mathcal F_i(x,y,t)
=\E_i\!\left[
x^{X_t}y^{Y_t}
\exp\!\left\{-\int_0^t
(\delta_1X_s+\delta_2Y_s)\,\dd s\right\}\right].
\end{equation*}
Their coefficients and total masses are
\begin{equation}
\label{eq:killed-coeff}
p_{m,n}^{(i)}(t)
:=[x^my^n]\mathcal F_i(x,y,t)
=\Prob_i\{X_t=m,\,Y_t=n,\,\tau_c>t\},
\end{equation}
\begin{equation}
\label{eq:sigma}
\sigma_i(t)
:=\mathcal F_i(1,1,t)
=
\begin{cases}
S(t), & i=1,\\
G(t), & i=2,
\end{cases}
\end{equation}
so that $1-\sigma_i(t)=\Prob_i\{\tau_c\le t\}$. Write
\begin{equation}
\label{eq:sigma-infty}
\sigma_i(\infty)
:=\lim_{t\to\infty}\sigma_i(t)
=
\begin{cases}
h, & i=1,\\
r_{2,-}, & i=2,
\end{cases}
\end{equation}
for the forever-avoidance probabilities of the main theorem.

The catastrophe-flux generating function of the appendix is
\begin{equation}
\label{eq:flux-pgf}
\mathcal J_i(x,y,t)
=
\delta_1 x\,\partial_x\mathcal F_i(x,y,t)
+
\delta_2 y\,\partial_y\mathcal F_i(x,y,t).
\end{equation}
Coefficientwise,
\begin{equation}
\label{eq:flux-coeff}
j_{m,n}^{(i)}(t)
:=[x^my^n]\mathcal J_i(x,y,t)
=
(\delta_1 m+\delta_2 n)\,p_{m,n}^{(i)}(t).
\end{equation}
From state $(m,n)$ the process is absorbed at rate $\delta_1m+\delta_2n$, so
$j_{m,n}^{(i)}(t)$ is the joint density (in $t$) of the event
\[
\bigl\{\tau_c\in\dd t,\;
X_{\tau_c-}=m,\;
Y_{\tau_c-}=n\bigr\}.
\]
In particular $j_{0,0}^{(i)}(t)\equiv 0$, since the empty state has catastrophe rate zero.
Summing over counts recovers the marginal catastrophe-time densities
\begin{equation}
\label{eq:flux-total}
\mathcal J_1(1,1,t)=-S'(t),
\qquad
\mathcal J_2(1,1,t)=-G'(t).
\end{equation}

\section{Joint distribution of the release counts}

\subsection{Point masses}

Integrating the flux density over all future times yields the law of the pre-catastrophe state on
the event that catastrophe eventually occurs. Adding the forever-avoidance atom at the origin
completes the law of $(\Xrel,\Yrel)$.

\begin{proposition}[Joint law of ultimate release]
\label{prop:joint-pmf}
For each initial type $i\in\{1,2\}$ and all integers $m,n\ge 0$,
\begin{equation}
\label{eq:joint-pmf}
\Prob_i\{\Xrel=m,\,\Yrel=n\}
=
\begin{cases}
\displaystyle
\sigma_i(\infty),
& (m,n)=(0,0),\\[1.1em]
\displaystyle
\int_0^\infty
(\delta_1 m+\delta_2 n)\,
p_{m,n}^{(i)}(t)\,
\dd t,
& m+n\ge 1.
\end{cases}
\end{equation}
Equivalently, for every $(m,n)\neq(0,0)$,
\begin{equation}
\label{eq:joint-pmf-j}
\Prob_i\{\Xrel=m,\,\Yrel=n\}
=
\int_0^\infty j_{m,n}^{(i)}(t)\,\dd t,
\end{equation}
while the origin mass is
\begin{equation}
\label{eq:origin-mass}
\Prob_i\{\Xrel=0,\,\Yrel=0\}
=
\Prob_i\{\tau_c=\infty\}
=
\sigma_i(\infty).
\end{equation}
\end{proposition}

\begin{proof}
On $\{\tau_c<\infty\}$ one has $(\Xrel,\Yrel)=(X_{\tau_c-},Y_{\tau_c-})$. The occupation measure of
the first entrance to $(\R,\R)$ from the pre-absorption state $(m,n)$ is exactly the integrated
flux \cref{eq:joint-pmf-j}. Since $j_{0,0}^{(i)}\equiv 0$, no catastrophe mass accumulates at the
origin. On $\{\tau_c=\infty\}$ the definition assigns $(0,0)$, which contributes the atom
$\sigma_i(\infty)$. These two pieces partition the sample space, so the displayed probabilities
sum to one:
\begin{align*}
\sum_{m,n\ge 0}
\Prob_i\{\Xrel=m,\,\Yrel=n\}
&=
\sigma_i(\infty)
+
\int_0^\infty
\sum_{m,n}
j_{m,n}^{(i)}(t)\,
\dd t
\\
&=
\sigma_i(\infty)
+
\int_0^\infty
\mathcal J_i(1,1,t)\,
\dd t
\\
&=
\sigma_i(\infty)
+
\int_0^\infty
\bigl(-\sigma_i'(t)\bigr)\,
\dd t
=
\sigma_i(\infty)
+
\bigl(1-\sigma_i(\infty)\bigr)
=1,
\end{align*}
where the last step uses $\sigma_i(0)=1$ and $\sigma_i\downarrow\sigma_i(\infty)$.
\end{proof}

\subsection{Joint probability generating function}

\begin{proposition}[Release PGF]
\label{prop:release-pgf}
The joint probability generating function of the ultimate release counts is
\begin{equation}
\label{eq:release-pgf}
\begin{aligned}
\mathcal P_i^{\mathrm{rel}}(x,y)
&:=
\E_i\bigl[x^{\Xrel}y^{\Yrel}\bigr]
\\
&=
\sigma_i(\infty)
+
\int_0^\infty
\mathcal J_i(x,y,t)\,
\dd t
\\
&=
\sigma_i(\infty)
+
\int_0^\infty
\bigl(
\delta_1 x\,\partial_x
+
\delta_2 y\,\partial_y
\bigr)
\mathcal F_i(x,y,t)\,
\dd t.
\end{aligned}
\end{equation}
Consequently the joint masses of \cref{prop:joint-pmf} are recovered by extraction,
\begin{equation}
\label{eq:extract}
\Prob_i\{\Xrel=m,\,\Yrel=n\}
=
[x^my^n]\,\mathcal P_i^{\mathrm{rel}}(x,y).
\end{equation}
\end{proposition}

\begin{proof}
Expand $\mathcal J_i$ by \cref{eq:flux-coeff} and integrate term by term. The constant term of the
integral vanishes, so the origin coefficient of $\mathcal P_i^{\mathrm{rel}}$ is exactly
$\sigma_i(\infty)$; every other coefficient is the integrated flux of \cref{eq:joint-pmf-j}.
\end{proof}

\begin{remark}[Finite horizon]
If one stops at a fixed observational horizon $T<\infty$ and records
\begin{equation*}
(\Xrel^{(T)},\Yrel^{(T)})
=
\begin{cases}
(X_{\tau_c-},Y_{\tau_c-}) & \text{on }\{\tau_c\le T\},\\
(0,0) & \text{on }\{\tau_c>T\},
\end{cases}
\end{equation*}
the same argument yields the truncated generating function
\begin{equation}
\label{eq:release-pgf-T}
\mathcal P_{i,T}^{\mathrm{rel}}(x,y)
=
\sigma_i(T)
+
\int_0^T
\mathcal J_i(x,y,t)\,
\dd t,
\end{equation}
with $\mathcal P_i^{\mathrm{rel}}=\lim_{T\to\infty}\mathcal P_{i,T}^{\mathrm{rel}}$ pointwise for
$|x|\le 1$, $|y|\le 1$.
\end{remark}

\section{Conditional laws and decomposition}

The appendix records the distribution of the pre-catastrophe state \emph{conditional on
catastrophe}. That object is the catastrophe part of \cref{eq:release-pgf}, renormalised.

\begin{proposition}[Catastrophe-conditional release]
\label{prop:conditional}
Assume $\sigma_i(\infty)<1$. The law of $(\Xrel,\Yrel)$ conditional on eventual catastrophe is
supported on $\{(m,n):m+n\ge 1\}$ and has probability generating function
\begin{equation}
\label{eq:R-cat}
\mathcal R_i^{\mathrm{cat}}(x,y)
:=
\E_i\bigl[
x^{\Xrel}y^{\Yrel}
\bigm|
\tau_c<\infty
\bigr]
=
\frac{
\displaystyle
\int_0^\infty
\mathcal J_i(x,y,t)\,
\dd t
}{
1-\sigma_i(\infty)
}.
\end{equation}
Equivalently, for $m+n\ge 1$,
\begin{equation}
\label{eq:conditional-pmf}
\Prob_i\{\Xrel=m,\,\Yrel=n\mid\tau_c<\infty\}
=
\frac{
\displaystyle
\int_0^\infty j_{m,n}^{(i)}(t)\,\dd t
}{
1-\sigma_i(\infty)
}.
\end{equation}
The unconditional release PGF therefore decomposes as the two-point mixture
\begin{equation}
\label{eq:mixture}
\mathcal P_i^{\mathrm{rel}}(x,y)
=
\sigma_i(\infty)
\cdot 1
+
\bigl(1-\sigma_i(\infty)\bigr)\,
\mathcal R_i^{\mathrm{cat}}(x,y),
\end{equation}
i.e.\ an atom of mass $\sigma_i(\infty)$ at $(0,0)$, and with complementary mass the
catastrophe-conditional release law of the appendix (the $T\to\infty$ limit of
$\mathcal R_{i,T}^{\mathrm{cat}}$).
\end{proposition}

\begin{proof}
On $\{\tau_c<\infty\}$ one has $(\Xrel,\Yrel)=(X_{\tau_c-},Y_{\tau_c-})$, so
\cref{eq:R-cat} is the integrated and renormalised flux already identified in the appendix.
\Cref{eq:mixture} is then immediate from \cref{eq:release-pgf}.
\end{proof}

Time-resolved versions are likewise available from the appendix: conditional on
$\tau_c=t$ in the density sense,
\begin{equation}
\label{eq:R-given-t}
\E_i\bigl[
x^{X_{\tau_c-}}y^{Y_{\tau_c-}}
\bigm|
\tau_c=t
\bigr]
=
\frac{\mathcal J_i(x,y,t)}{\mathcal J_i(1,1,t)},
\end{equation}
and therefore
\begin{equation}
\label{eq:pmf-given-t}
\Prob_i\{X_{\tau_c-}=m,\,Y_{\tau_c-}=n\mid\tau_c=t\}
=
\frac{j_{m,n}^{(i)}(t)}{\mathcal J_i(1,1,t)}.
\end{equation}
Integrating \cref{eq:R-given-t} against the catastrophe-time density
$\mathcal J_i(1,1,t)=-\sigma_i'(t)$ recovers the numerator of \cref{eq:R-cat}.

\section{Marginal distributions}

The marginal generating functions are the obvious specialisations of
$\mathcal P_i^{\mathrm{rel}}$.

\begin{corollary}[Marginals of release]
\label{cor:marginals}
For each initial type $i$,
\begin{align}
\E_i\bigl[x^{\Xrel}\bigr]
&=
\mathcal P_i^{\mathrm{rel}}(x,1)
=
\sigma_i(\infty)
+
\int_0^\infty
\delta_1 x\,\partial_x\mathcal F_i(x,1,t)\,
\dd t
+
\int_0^\infty
\delta_2\,\partial_y\mathcal F_i(x,1,t)\,
\dd t,
\label{eq:marg-X}
\\
\E_i\bigl[y^{\Yrel}\bigr]
&=
\mathcal P_i^{\mathrm{rel}}(1,y)
=
\sigma_i(\infty)
+
\int_0^\infty
\delta_1\,\partial_x\mathcal F_i(1,y,t)\,
\dd t
+
\int_0^\infty
\delta_2 y\,\partial_y\mathcal F_i(1,y,t)\,
\dd t.
\label{eq:marg-Y}
\end{align}
The corresponding point masses are
\begin{align}
\Prob_i\{\Xrel=m\}
&=
[x^m]\,\mathcal P_i^{\mathrm{rel}}(x,1),
\label{eq:marg-X-pmf}
\\
\Prob_i\{\Yrel=n\}
&=
[y^n]\,\mathcal P_i^{\mathrm{rel}}(1,y).
\label{eq:marg-Y-pmf}
\end{align}
In particular both marginals carry the common origin atom $\sigma_i(\infty)$, since
$(\Xrel,\Yrel)=(0,0)$ on $\{\tau_c=\infty\}$.
Explicitly, for $m\ge 1$,
\begin{equation}
\label{eq:marg-X-explicit}
\Prob_i\{\Xrel=m\}
=
\int_0^\infty
\sum_{n=0}^\infty
(\delta_1 m+\delta_2 n)\,
p_{m,n}^{(i)}(t)\,
\dd t,
\end{equation}
and for $n\ge 1$,
\begin{equation}
\label{eq:marg-Y-explicit}
\Prob_i\{\Yrel=n\}
=
\int_0^\infty
\sum_{m=0}^\infty
(\delta_1 m+\delta_2 n)\,
p_{m,n}^{(i)}(t)\,
\dd t,
\end{equation}
while
\begin{equation}
\label{eq:marg-origin}
\Prob_i\{\Xrel=0\}
=
\sigma_i(\infty)
+
\int_0^\infty
\sum_{n=1}^\infty
\delta_2 n\,
p_{0,n}^{(i)}(t)\,
\dd t,
\end{equation}
and likewise
\begin{equation}
\label{eq:marg-Y-zero}
\Prob_i\{\Yrel=0\}
=
\sigma_i(\infty)
+
\int_0^\infty
\sum_{m=1}^\infty
\delta_1 m\,
p_{m,0}^{(i)}(t)\,
\dd t.
\end{equation}
(The second summands are the probabilities of catastrophe from a pure type-$2$ or pure type-$1$
state, respectively; they vanish if the corresponding catastrophe rate is zero.)
\end{corollary}

\begin{remark}[Joint origin versus marginal origins]
The event $\{\Xrel=0\}$ is strictly larger than $\{\Xrel=\Yrel=0\}$ whenever catastrophe can
occur from a state with $X=0$ and $Y\ge 1$. Thus
$\Prob_i\{\Xrel=0\}\ge\sigma_i(\infty)$, with equality if and only if pure type-$2$ states cannot
trigger (in particular if $\delta_2=0$). The same remark applies to $\Yrel$ with the roles of the
types reversed.
\end{remark}

\section{Means and low-order moments}

Differentiating the release PGF at $(1,1)$ yields moments of the ultimate release.
Because $\mathcal P_i^{\mathrm{rel}}$ has an atom at the origin, these moments are
\emph{unconditional} means of released load (including zeros on forever-avoidance paths).
Conditional means given catastrophe are obtained by dividing by $1-\sigma_i(\infty)$.

\begin{corollary}[Mean released load]
\label{cor:means}
Whenever the derivatives exist in a neighbourhood of $(1,1)$,
\begin{align}
\E_i[\Xrel]
&=
\partial_x\mathcal P_i^{\mathrm{rel}}(1,1)
=
\int_0^\infty
\Bigl(
\delta_1\,
\partial_x\bigl(x\,\partial_x\mathcal F_i\bigr)
+
\delta_2\,
\partial_x\bigl(y\,\partial_y\mathcal F_i\bigr)
\Bigr)
\Big|_{(x,y)=(1,1)}
\dd t,
\label{eq:mean-X}
\\
\E_i[\Yrel]
&=
\partial_y\mathcal P_i^{\mathrm{rel}}(1,1)
=
\int_0^\infty
\Bigl(
\delta_1\,
\partial_y\bigl(x\,\partial_x\mathcal F_i\bigr)
+
\delta_2\,
\partial_y\bigl(y\,\partial_y\mathcal F_i\bigr)
\Bigr)
\Big|_{(x,y)=(1,1)}
\dd t.
\label{eq:mean-Y}
\end{align}
Equivalently, in coefficient form,
\begin{align}
\E_i[\Xrel]
&=
\int_0^\infty
\sum_{m,n\ge 0}
m\,(\delta_1 m+\delta_2 n)\,
p_{m,n}^{(i)}(t)\,
\dd t,
\label{eq:mean-X-sum}
\\
\E_i[\Yrel]
&=
\int_0^\infty
\sum_{m,n\ge 0}
n\,(\delta_1 m+\delta_2 n)\,
p_{m,n}^{(i)}(t)\,
\dd t.
\label{eq:mean-Y-sum}
\end{align}
The conditional means given eventual catastrophe are
\begin{equation}
\label{eq:cond-means}
\E_i[\Xrel\mid\tau_c<\infty]
=
\frac{\E_i[\Xrel]}{1-\sigma_i(\infty)},
\qquad
\E_i[\Yrel\mid\tau_c<\infty]
=
\frac{\E_i[\Yrel]}{1-\sigma_i(\infty)}.
\end{equation}
\end{corollary}

(The origin atom does not contribute to either mean.)

\section{Evaluation from the closed killed PGF}

The formulae above are distributional identities: they do not require a closed form for
$\mathcal F_i$. When the hypergeometric solution of the appendix is available, they become
explicit integral transforms of that solution.

In the generic regime of the main theorem, after the type-$1$ time scaling
$\hat t=\lambda_1 t$,
\begin{align}
\widehat{\mathcal F}_2(x,y,\hat t)
&=
r_{2,-}-\alpha_2
\frac{z_2(\hat t;y)}{1-z_2(\hat t;y)},
\label{eq:F2-closed}
\\
\widehat{\mathcal F}_1(x,y,\hat t)
&=
h-\hat q_2 z
\frac{K_1W_1'(z)+K_2W_2'(z)}
{K_1W_1(z)+K_2W_2(z)},
\quad
z=z_2(\hat t;y),
\label{eq:F1-closed}
\end{align}
with $z_{2,0}(y)=(y-r_{2,-})/(y-r_{2,+})$,
$z_2(\hat t;y)=z_{2,0}(y)\,e^{\hat q_2\hat t}$,
and Wronskian coefficients $(K_1,K_2)$ built from the initial logarithmic derivative
$L(x,y)=(h-x)/(\hat q_2 z_{2,0}(y))$ as in the appendix.
Substituting \cref{eq:F1-closed,eq:F2-closed} into \cref{eq:flux-pgf} and integrating in physical
time (or in $\hat t$, with the Jacobian $1/\lambda_1$) yields
$\mathcal P_i^{\mathrm{rel}}$ as an explicit improper integral of elementary and hypergeometric
functions of $(x,y)$. Coefficients may then be read by a bivariate Cauchy integral on a product
of small circles about the origin, or by series expansion of $\mathcal F_i$ before the time
integral when a power-series route is preferred.

For numerical work at fixed $(x,y)$, one may equally integrate the killed backward system of the
appendix for $\mathcal F_i(x,y,\cdot)$ and accumulate the flux integrand of \cref{eq:release-pgf}
in parallel, without extracting coefficients first. The joint masses
\cref{eq:joint-pmf} are the special case of integrating the scalar ODEs for each coefficient
$p_{m,n}^{(i)}(t)$ against the weight $\delta_1m+\delta_2n$.

\section{Summary}

The ultimate release pair $(\Xrel,\Yrel)$ is completely determined by the killed generating
function through its catastrophe flux:
\begin{equation}
\label{eq:summary}
\boxed{
\begin{aligned}
\mathcal P_i^{\mathrm{rel}}(x,y)
&=
\E_i\bigl[x^{\Xrel}y^{\Yrel}\bigr]
=
\sigma_i(\infty)
+
\int_0^\infty
\mathcal J_i(x,y,t)\,
\dd t,
\\[0.4em]
\Prob_i\{\Xrel=m,\,\Yrel=n\}
&=
[x^my^n]\,\mathcal P_i^{\mathrm{rel}}(x,y)
\\
&=
\begin{cases}
\sigma_i(\infty),
& (m,n)=(0,0),\\[0.35em]
\displaystyle
\int_0^\infty
(\delta_1 m+\delta_2 n)\,
p_{m,n}^{(i)}(t)\,
\dd t,
& m+n\ge 1.
\end{cases}
\end{aligned}
}
\end{equation}
The catastrophe-conditional law of the appendix is the same object with the origin atom removed
and the remaining mass renormalised by $1-\sigma_i(\infty)$. Marginals and moments follow by
specialising or differentiating $\mathcal P_i^{\mathrm{rel}}$.

\end{document}
